\documentclass{notaaula}

        \titulonota{Trabalho e energia mecânica}
        \creditos{Turma 1C · Física · Outubro/2026 · Prof. Josué Cavalcante}

        \begin{document}
        \cabecalho

        \begin{sobre}
        Dinâmica: energia, trabalho de uma força, formas de energia e conservação da energia mecânica. Habilidades em foco: EM13CNT207, EM13CNT301, EM13CNT306 e EM13CNT307.
        \end{sobre}

        \section{Trabalho e transferência de energia}

\begin{copiar}{Trabalho de uma força constante}
O trabalho mede a energia transferida por uma força durante um deslocamento:
\[
  W=F\,d\cos\theta.
\]
Ele é positivo quando a força favorece o movimento, negativo quando se opõe e nulo quando a
força é perpendicular ao deslocamento. A unidade é o joule: $1\un{J}=1\un{N\,m}$.
\simbolos{$W$ (J); $F$ (N); $d$ (m); $\theta$ é o ângulo entre força e deslocamento.}
\end{copiar}

\begin{exemplo}
Uma pessoa empurra uma caixa com força constante de $50\un{N}$, paralela ao piso, por $4\un{m}$.
\[
  W=50\cdot4\cdot\cos0^\circ=\resultado{200\un{J}}.
\]
Se o atrito vale $20\un{N}$, seu trabalho é $W_{at}=-20\cdot4=-80\un{J}$.
\end{exemplo}

\section{Energia cinética e potenciais}

\begin{copiar}{Formas mecânicas}
\formulas
\[
  E_c=\frac{mv^2}{2},\qquad
  E_{pg}=mgh,\qquad
  E_{pe}=\frac{kx^2}{2}.
\]
A energia cinética depende da rapidez. A energia potencial gravitacional depende do nível de
referência escolhido. A energia potencial elástica aparece na deformação de molas ideais.
\relacoes O trabalho da força resultante é a variação da energia cinética:
\[
  W_{\mathrm{res}}=\Delta E_c.
\]
\end{copiar}

\begin{dica}
Dobrar a rapidez quadruplica $E_c$; dobrar a altura duplica $E_{pg}$. Observe o expoente antes
de usar uma regra de proporção.
\end{dica}

\section{Conservação da energia mecânica}

\begin{copiar}{Sistema sem dissipação}
A energia mecânica é
\[
  E_m=E_c+E_p.
\]
Quando apenas forças conservativas realizam trabalho, $E_{m,i}=E_{m,f}$. Com atrito ou outra
dissipação, parte da energia mecânica se transforma em energia interna, som ou deformação:
\[
  E_{m,i}+W_{\mathrm{nc}}=E_{m,f}.
\]
A energia total continua conservada; o que pode diminuir é a parcela mecânica.
\diaadia{Freios transformam energia cinética principalmente em energia interna nos discos e pastilhas.}
\end{copiar}

\begin{exemplo}[Queda sem resistência do ar]
Uma pedra parte do repouso de $5\un{m}$ de altura. Com $g=10\un{m/s^2}$,
\[
  mgh=\frac{mv^2}{2}
  \Rightarrow v=\sqrt{2gh}=\sqrt{100}
  =\resultado{10\un{m/s}}.
\]
A massa cancela porque todos os corpos, sem resistência do ar, têm a mesma aceleração gravitacional.
\end{exemplo}

\section{Potência e rendimento}

\begin{copiar}{Rapidez da transformação}
\[
  \mathcal P=\frac{W}{\Delta t}=\frac{\Delta E}{\Delta t},
  \qquad
  \eta=\frac{E_{\mathrm{útil}}}{E_{\mathrm{total}}}.
\]
Potência é medida em watt: $1\un{W}=1\un{J/s}$. O rendimento real fica entre 0 e 1, ou entre
$0\%$ e $100\%$.
\end{copiar}

\begin{exercicios}
\nivel{1}{Conceitos}
\begin{questoes}
\item Quando o trabalho de uma força é nulo por causa do ângulo?
\item Uma força de $30\un{N}$ atua no sentido de um deslocamento de $5\un{m}$. Calcule $W$.
\item Calcule $E_c$ de um corpo de $2\un{kg}$ a $3\un{m/s}$.
\item Calcule $E_{pg}$ de $4\un{kg}$ a $6\un{m}$, com $g=10\un{m/s^2}$.
\item O que acontece com $E_c$ quando a rapidez dobra?
\nivel{2}{Aplicação}
\item Um motor realiza $6000\un{J}$ em $20\un{s}$. Qual é sua potência média?
\item Um equipamento recebe $500\un{J}$ e fornece $400\un{J}$ úteis. Ache o rendimento.
\item Um corpo cai do repouso de $20\un{m}$. Despreze o ar e determine a rapidez ao chegar ao solo.
\item Uma mola de $k=200\un{N/m}$ é comprimida $\dec{0,10}\un{m}$. Calcule $E_{pe}$.
\nivel{3}{Síntese}
\item Uma caixa recebe trabalho de $120\un{J}$ da força aplicada e $-40\un{J}$ do atrito. Qual é $\Delta E_c$?
\item Explique por que energia mecânica e energia total não são sinônimos.
\item Uma montanha-russa perde $30\%$ da energia mecânica inicial de $20\,000\un{J}$. Quanto resta?
\end{questoes}
\gabarito{1) $90^\circ$; 2) $150\un{J}$; 3) $9\un{J}$; 4) $240\un{J}$;
5) quadruplica; 6) $300\un{W}$; 7) $80\%$; 8) $20\un{m/s}$;
9) $1\un{J}$; 10) $80\un{J}$; 12) $14\,000\un{J}$.}
\end{exercicios}


\end{document}
